\section{Introduction}

Modern systems increasingly construct task-specific filesystem views around
ordinary tools. Agents branch, fork, checkpoint, hide side effects, and run
build or test commands over repository
workspaces~\cite{agentfs_repo,branchfs_repo,yolofs_repo}. Environment systems
validate repositories against generated Docker state, cache epochs, and verified
local artifacts~\cite{docker_overview,swe_factory_repo,menvagent_repo,swe_rebench_v2_repo},
while services consume projected configuration, secrets, certificates, or
fixtures as operational state changes~\cite{kubernetes_projected_volumes,kubernetes_configmaps,kubernetes_secrets}.

However, these views make pathname binding a correctness boundary. Binding a
name to the wrong lower object can validate the wrong revision, leak a write
into a parent workspace, feed stale cache input into a build, or expose the
wrong service configuration. In many of these cases, the needed behavior is not
a new file format, storage service, or custom file operation. The needed
behavior chooses which existing object a name denotes, or whether that object is
visible.

The root cause is that existing mechanisms couple pathname binding to a broader
owner. Namespace construction binds names before use by materializing or
updating mounts, layers, links, or files. FUSE, stackable, and custom
filesystems can instead choose a binding during lookup, but that choice lives
inside a filesystem implementation that also owns a wider operation and
lifetime boundary. Workload state may need to change only which existing object
a name denotes, while data, writes, permissions, page-cache behavior,
persistence, and consistency should remain with the lower filesystem.

Existing mechanisms bracket this lookup-time selection problem. Namespace
mechanisms such as bind mounts, OverlayFS, projected volumes, and copied trees
construct or materialize a namespace before the application uses
it~\cite{linux_mount,linux_overlayfs,kubernetes_projected_volumes}.
eBPF LSM hooks can attach verified access-control policy to security
decisions~\cite{linux_bpf_lsm},
but they are not an object-selection boundary. FUSE, stackable filesystems,
custom filesystems, and metadata services can compute richer views, but they
cross into a broader filesystem boundary~\cite{linux_fuse,wrapfs,bento,deltafs_sc21,indexfs,tablefs}.
They may be the right mechanism when the oracle needs synthetic contents,
custom metadata, data-path mediation, or runtime orchestration, but they are
broader than lookup-time name-resolution policy.

We argue that dynamic filesystem views are a pathname late-binding problem,
not necessarily a new-filesystem problem. The recurring subproblem is a
state-dependent path view: workload state changes which existing lower object
a pathname denotes or whether that object is visible. VFS name resolution is
where this binding is chosen. The key separation is therefore to make binding
policy programmable while the VFS and lower filesystem retain object ownership
and ordinary filesystem semantics.

The central challenge is defining when the VFS can accept a policy-selected
object as a normal name-resolution result. An in-progress path walk is not an
inode pointer: \code{nameidata} carries operation intent, the current root and
mount, traversal restrictions, RCU/ref-walk state, and sequence numbers used to
validate concurrent rename and mount activity. Replacing only its path can
bypass those constraints or leave the walk with inconsistent ownership. The
extension must therefore admit only pre-authorized, type- and
operation-compatible objects; obtain stable lifetime without exposing kernel
pointers to BPF; reject selections whose path constraints cannot be preserved;
and then resume the original VFS state machine rather than emulate it.

Two additional requirements follow from where this object-selection boundary
sits. Intermediate components, final open/create components, and directory
entries traverse different VFS paths, but must implement the same binding
contract. Cache-hot lookup performs that work during lockless RCU-walk, while
target registration and replacement occur concurrently. Finally, the decision
point lies on a system-wide hot path, so disabled and unrelated lookups must
bypass policy work before context construction.

We present \namei, which separates pathname-binding policy from filesystem
ownership. One eBPF function returns bounded pass, redirect, hide, or
\code{SELECT_TARGET} actions at the VFS paths that observe names. Opaque target
IDs refer to kernel-held \code{struct path} objects; the kernel validates each
result, converts an RCU-borrowed target into stable namei state, and resumes
ordinary permission, open, exec, and lower-filesystem operations on the
selected object. A static key and an RCU parent-scope filter bypass policy work
before context construction when no relevant policy exists. The resulting
ownership split is analogous to \code{sched_ext}: policy is programmable,
while the kernel retains the subsystem machinery that owns objects and
executes semantics~\cite{linux_sched_ext}.

We evaluate whether this ownership split is sufficient for source-derived
path-view policies, cheaper than feature-equivalent FUSE placement, and
narrower in method and runtime responsibility than broader filesystem
ownership. RQ1 asks whether KVM workloads pass
the same correctness conditions while preserving lower-filesystem semantics.
Reviewed formal KVM runs pass fixed oracles for Agent workspaces, official XDG
Documents portal application-sharing grants, concurrent Bazel action views, the
already-materialized payload-view subset of Kubernetes \code{AtomicWriter},
and toolchain environment selection. The headline Agent result includes three
fresh boots of
a released Click source task: concurrent base and completed views produce
39/40 and 40/40 tests, and switch/rollback changes both the selected object and
the task outcome. RQ2 compares
the same policies with feature-equivalent FUSE: the matched Agent lifecycle is
11.32\(\times\) [11.24, 11.64], and cache-hot FxMark \code{SELECT} reaches
1.052--1.088\(\times\) cached-FUSE throughput across MRPL, MRPM, and MRPH.
On corrected FxMark directory enumeration, \code{SELECT} reaches
2.20--3.66\(\times\) FUSE throughput in five of six cells; at four
shared-directory workers, the ratio is 1.018 [0.907, 1.135], exposing the
shared-directory contention boundary.
RQ3 executes a matched
Wrapfs-derived implementation; both mechanisms pass 37/37 workspace oracles
and all 21 fault cells in each of three boots, while runtime tracing exposes
the broader stackable filesystem-method boundary.

The paper's primary contribution is the design and Linux implementation of a
late-binding boundary: programmable selection of existing objects during VFS
name resolution while preserving VFS and lower-filesystem ownership.
The paper makes two concrete contributions:
\begin{itemize}
\item the design and modified-Linux implementation of a pathname-binding
extension, including one contract across component lookup, final lookup, and
readdir, registered-target lifetime across RCU/ref-walk, preservation of normal
VFS processing after selection, and a scope-aware fast path
(Sections~\ref{sec:design}
and~\ref{sec:implementation}); and
\item an evaluation that holds workload correctness fixed while comparing
\namei with feature-equivalent FUSE and recording which responsibilities would
move into custom or stackable filesystems (Section~\ref{sec:evaluation}).
\end{itemize}
